\section{When temporary overparameterization can pay for itself}
The previous compute-to-accuracy result asks whether a contraction is beneficial relative to continuing the same wide state. A harder question is the one relevant for the claim that a network should \emph{learn wide and then become small}: why not train the final compact architecture from the beginning? The answer cannot be ``because contraction is safe.'' Safety only says that little is lost at the contraction event. Temporary width is useful at equal compute only if the wide phase creates a sufficiently better compact state to repay the compute spent while wide.

Fix a wide architecture $W$ and a realizable compact architecture $r$, with per-update costs $C_W\ge C_r>0$. Starting from a common target-free initialization rule, let
\[
R_{W\to r}(t,H)
\]
denote the risk obtained by training the wide model for $t$ updates, structurally contracting it to architecture $r$, and then training the compact model for another $H$ updates. Let $R_r(k)$ denote the risk of the compact-from-start control after $k$ updates. Define the equal-compute compact budget
\[
N_r(t,H)
:=
\left\lceil\frac{C_Wt+C_rH}{C_r}\right\rceil,
\qquad
K(t,H):=t+H.
\]
The ceiling weakly favors the compact control and can be dropped when the compute ratio is exact.

\begin{definition}[Discovery advantage and compute opportunity gain]
Define
\[
D_r(t,H)
:=R_r(K(t,H))-R_{W\to r}(t,H)
\]
and
\[
O_r(t,H)
:=R_r(K(t,H))-R_r(N_r(t,H)).
\]
$D_r$ measures how much better the state produced by temporary width is than the compact-from-start state at the \emph{same update count}. It includes every useful effect of temporary overparameterization---representation discovery, optimization-path selection, conditioning, and initialization amplification---after both paths have been expressed in the same final architecture. $O_r$ measures how much the compact control gains from the additional updates it can buy with the compute that the wide phase consumed.
\end{definition}

\begin{theorem}[Exact break-even identity]\label{thm:temporary-width-break-even}
For every $t,H\ge0$,
\[
\boxed{
R_{W\to r}(t,H)-R_r(N_r(t,H))
=
O_r(t,H)-D_r(t,H).}
\]
Consequently, temporary overparameterization wins the equal-compute comparison if and only if
\[
\boxed{D_r(t,H)>O_r(t,H).}
\]
If $R_r(k)$ is nonincreasing in $k$ and $C_W\ge C_r$, then $O_r(t,H)\ge0$. In that case $D_r(t,H)>0$ is a necessary condition for any equal-compute advantage.
\end{theorem}

\begin{proof}
Add and subtract $R_r(K(t,H))$:
\begin{align*}
R_{W\to r}(t,H)-R_r(N_r(t,H))
&=
\bigl[R_{W\to r}(t,H)-R_r(K(t,H))\bigr]\\
&\quad+
\bigl[R_r(K(t,H))-R_r(N_r(t,H))\bigr]\\
&=-D_r(t,H)+O_r(t,H).
\end{align*}
The first equivalence follows immediately. Moreover $C_W\ge C_r$ implies $N_r(t,H)\ge K(t,H)$; monotonicity of the compact risk then gives $R_r(K)\ge R_r(N_r)$ and hence $O_r\ge0$.
\end{proof}

\begin{corollary}[Safe contraction is not enough]\label{cor:safety-not-enough}
Suppose the structural contraction at time $t$ has arbitrarily small finite-horizon damage relative to continuing the wide model. This fact alone places no sign restriction on $D_r(t,H)$ and therefore cannot imply an equal-compute advantage over compact-from-start. In particular, if $D_r(t,H)\le0$ and compact risk is nonincreasing, temporary overparameterization cannot win at equal compute.
\end{corollary}

\begin{proof}
The contraction-damage comparison uses a continuation of the checkpointed wide state as its reference, whereas $D_r$ compares against an independently trained compact path. These are different counterfactuals. The final statement follows from Theorem~\ref{thm:temporary-width-break-even} because $O_r\ge0$.
\end{proof}

\begin{remark}[The empirical decomposition that should be reported]
An equal-compute experiment should report all three quantities
\[
R_{W\to r}(t,H)-R_r(K),
\qquad
O_r(t,H),
\qquad
R_{W\to r}(t,H)-R_r(N_r),
\]
rather than only the last one. The first term reveals whether temporary width actually discovered a better compact state; the second is the price of having trained wide; the third is their net effect. This decomposition distinguishes a failed discovery mechanism from a useful discovery mechanism that is simply too expensive.
\end{remark}

\begin{remark}[Research consequence]
The break-even theorem changes the role of the controller. A safe-pruning certificate determines when one \emph{may} contract. A useful learn-wide-then-shrink algorithm needs an additional economic gate: estimated discovery advantage must be large enough to have justified the historical wide-phase compute, or, prospectively, the expected future discovery value of remaining wide must exceed its incremental compute price. This is stronger than spectral truncation and stronger than phase separation alone.
\end{remark}
